\chapter{8-3编码器}

\section{第一步：解决什么问题}

想象一个场景：你的FPGA有8个按键，编号0到7。用户按下其中一个按键。

\begin{itemize}
  \item 如果用8根信号线传给处理器，需要8个引脚——浪费资源
  \item 更好的方案：用3根信号线传递"哪个按键被按下"的信息
  \item 因为 $2^3 = 8$，3根线恰好能表示8种状态
\end{itemize}

\textbf{编码器（Encoder）的本质}：将 $2^n$ 个输入信号，压缩成 $n$ 位二进制输出。

\begin{keypoint}[编码器的本质]
编码器 = 信号压缩器：$2^n$ 个输入 $\to$ $n$ 位输出

8-3编码器：8个输入 $\to$ 3位输出
\end{keypoint}

\section{第二步：输入输出关系}

\begin{itemize}
  \item \textbf{输入}：$I_0, I_1, I_2, I_3, I_4, I_5, I_6, I_7$（8个输入信号，高电平有效）
  \item \textbf{输出}：$Y_2, Y_1, Y_0$（3位二进制编码输出）
  \item \textbf{约束}：任意时刻最多只有一个输入为1
\end{itemize}

当 $I_k = 1$（且其他输入为0）时，输出 $(Y_2 Y_1 Y_0)$ 应该等于 $k$ 的二进制表示。

\section{第三步：真值表}

\begin{center}
\begin{tabular}{|c|c|c|c|}
\hline
\textbf{有效输入} & $Y_2$ & $Y_1$ & $Y_0$ \\
\hline
$I_0 = 1$ & 0 & 0 & 0 \\
$I_1 = 1$ & 0 & 0 & 1 \\
$I_2 = 1$ & 0 & 1 & 0 \\
$I_3 = 1$ & 0 & 1 & 1 \\
$I_4 = 1$ & 1 & 0 & 0 \\
$I_5 = 1$ & 1 & 0 & 1 \\
$I_6 = 1$ & 1 & 1 & 0 \\
$I_7 = 1$ & 1 & 1 & 1 \\
\hline
\end{tabular}
\end{center}

\section{第四步：逻辑推导}

观察真值表，找出每个输出位什么时候为1：

\begin{itemize}
  \item $Y_2 = 1$ 当 $I_4=1$ 或 $I_5=1$ 或 $I_6=1$ 或 $I_7=1$
  \item $Y_1 = 1$ 当 $I_2=1$ 或 $I_3=1$ 或 $I_6=1$ 或 $I_7=1$
  \item $Y_0 = 1$ 当 $I_1=1$ 或 $I_3=1$ 或 $I_5=1$ 或 $I_7=1$
\end{itemize}

因此，逻辑表达式为：

\begin{align}
  Y_2 &= I_4 + I_5 + I_6 + I_7 \\
  Y_1 &= I_2 + I_3 + I_6 + I_7 \\
  Y_0 &= I_1 + I_3 + I_5 + I_7
\end{align}

每个输出位都是若干输入的\textbf{逻辑或（OR）}。

\section{第五步：结构化设计}

8-3编码器内部结构：

\begin{center}
\begin{tikzpicture}
  % Inputs
  \foreach \x in {0,...,7} {
    \node (I\x) at (0, -1.2*\x) {$I_\x$};
  }
  % OR gates
  \node[american or port,rotate=-90] (or2) at (4,-2.5) {};
  \node[american or port,rotate=-90] (or1) at (4,-5) {};
  \node[american or port,rotate=-90] (or0) at (4,-7.5) {};

  % Outputs
  \node (Y2) at (5.5,-2.5) {$Y_2$};
  \node (Y1) at (5.5,-5) {$Y_1$};
  \node (Y0) at (5.5,-7.5) {$Y_0$};

  % Connections for Y2
  \draw (I4) -- ++(1.5,0) |- (or2.in 1);
  \draw (I5) -- ++(1.5,0) |- (or2.in 2);
  \draw (I6) -- ++(2.0,0) |- (or2.in 1);
  \draw (I7) -- ++(2.5,0) |- (or2.in 2);

  \draw (or2.out) -- (Y2);
\end{tikzpicture}
\end{center}

\textbf{设计规律}：$Y_k$ 的输出 = 所有编号第 $k$ 位为1的输入的OR。

\begin{itemize}
  \item $Y_2$（bit 2 = 4的位）：连接 $I_4, I_5, I_6, I_7$（这些编号的二进制第2位都是1）
  \item $Y_1$（bit 1 = 2的位）：连接 $I_2, I_3, I_6, I_7$
  \item $Y_0$（bit 0 = 1的位）：连接 $I_1, I_3, I_5, I_7$
\end{itemize}

\section{第六步：为什么这样设计}

\begin{enumerate}
  \item \textbf{为什么用OR门？}

  观察真值表：对于 $Y_k$，任何使 $Y_k=1$ 的输入条件之间是"或"的关系——只要其中任何一个条件满足，$Y_k$ 就是1。

  \item \textbf{为什么输入之间存在约束？}

  如果两个输入同时为1，OR门的输出仍然是1，但编码结果就错了（对应了另一个不需要的编码）。所以编码器假设同一时刻只有一个输入有效。

  \item \textbf{优先级编码器的改进动机？}

  当多个输入同时有效时，普通编码器产生错误输出。优先级编码器增加了优先级逻辑——一般编号大的输入优先级高。这需要用if-else级联逻辑而非简单OR门。
\end{enumerate}

\section{第七步：考试常见变式}

\begin{enumerate}
  \item \textbf{变式1：16-4编码器}

  将16个输入编码为4位输出。原理完全相同，只是规模扩大。$Y_3 = I_8 + I_9 + \cdots + I_{15}$，以此类推。

  \item \textbf{变式2：优先级编码器（Priority Encoder）}

  允许同时多个输入为1，输出最高优先级输入的编号。这是考试的热门考点。

  优先级编码器的真值表（优先级：$I_7 > I_6 > \cdots > I_0$）：

  \begin{center}
  \begin{tabular}{|c|c|c|c|c|}
  \hline
  $I_7$ & $I_6$ & $I_5$ & ... & $(Y_2 Y_1 Y_0)$ \\
  \hline
  1 & x & x & ... & 111 \\
  0 & 1 & x & ... & 110 \\
  0 & 0 & 1 & ... & 101 \\
  \vdots & \vdots & \vdots & \ddots & \vdots \\
  \hline
  \end{tabular}
  \end{center}

  \item \textbf{变式3：带使能端的编码器}

  增加一个Enable输入，E=1时编码器正常工作，E=0时所有输出为0。

  \item \textbf{变式4：输出有效的编码器}

  增加一个GS（Group Select）输出，当任何输入有效时GS=1，否则GS=0。用于区分"选择了$I_0$"（输出000）和"没有选择"（输出也是000）。

  \item \textbf{变式5：低电平有效的编码器}

  输入和输出都是低电平有效（即0表示有效，1表示无效）。逻辑门类型可能需要相应改变。
\end{enumerate}

\section{第八步：VHDL实现}

\subsection{普通8-3编码器}

\begin{lstlisting}[caption={8-3编码器VHDL实现（if-else版本）}]
library ieee;
use ieee.std_logic_1164.all;

entity encoder_8to3 is
  port (
    I : in  std_logic_vector(7 downto 0);
    Y : out std_logic_vector(2 downto 0)
  );
end encoder_8to3;

architecture behavioral of encoder_8to3 is
begin
  process(I)
  begin
    case I is
      when "00000001" => Y <= "000";
      when "00000010" => Y <= "001";
      when "00000100" => Y <= "010";
      when "00001000" => Y <= "011";
      when "00010000" => Y <= "100";
      when "00100000" => Y <= "101";
      when "01000000" => Y <= "110";
      when "10000000" => Y <= "111";
      when others      => Y <= "000";
    end case;
  end process;
end behavioral;
\end{lstlisting}

\subsection{优先级编码器}

\begin{lstlisting}[caption={8-3优先级编码器（使用if-else级联）}]
library ieee;
use ieee.std_logic_1164.all;

entity priority_encoder_8to3 is
  port (
    I : in  std_logic_vector(7 downto 0);
    Y : out std_logic_vector(2 downto 0);
    GS: out std_logic
  );
end priority_encoder_8to3;

architecture behavioral of priority_encoder_8to3 is
begin
  process(I)
  begin
    GS <= '1';
    if    I(7) = '1' then Y <= "111";
    elsif I(6) = '1' then Y <= "110";
    elsif I(5) = '1' then Y <= "101";
    elsif I(4) = '1' then Y <= "100";
    elsif I(3) = '1' then Y <= "011";
    elsif I(2) = '1' then Y <= "010";
    elsif I(1) = '1' then Y <= "001";
    elsif I(0) = '1' then Y <= "000";
    else
      Y  <= "000";
      GS <= '0';  -- 无有效输入
    end if;
  end process;
end behavioral;
\end{lstlisting}

\section{第九步：Testbench验证}

\begin{lstlisting}[caption={8-3编码器Testbench}]
library ieee;
use ieee.std_logic_1164.all;

entity encoder_8to3_tb is
end encoder_8to3_tb;

architecture test of encoder_8to3_tb is
  signal I : std_logic_vector(7 downto 0);
  signal Y : std_logic_vector(2 downto 0);

  component encoder_8to3 is
    port (
      I : in  std_logic_vector(7 downto 0);
      Y : out std_logic_vector(2 downto 0)
    );
  end component;
begin
  uut: encoder_8to3 port map (I, Y);

  stimulus: process
  begin
    -- 测试每个输入
    I <= "00000001"; wait for 10 ns;
    I <= "00000010"; wait for 10 ns;
    I <= "00000100"; wait for 10 ns;
    I <= "00001000"; wait for 10 ns;
    I <= "00010000"; wait for 10 ns;
    I <= "00100000"; wait for 10 ns;
    I <= "01000000"; wait for 10 ns;
    I <= "10000000"; wait for 10 ns;
    -- 无输入
    I <= "00000000"; wait for 10 ns;
    wait;
  end process;
end test;
\end{lstlisting}

\textbf{验证要点}：
\begin{itemize}
  \item $I_0$ 有效时，$Y = 000$
  \item $I_7$ 有效时，$Y = 111$
  \item 中间值按二进制递增
\end{itemize}

\section{第十步：如何快速解题}

\begin{examtip}[编码器快速解题法]
\begin{enumerate}
  \item \textbf{确定输入输出数量}：$2^n$ 输入对应 $n$ 位输出
  \item \textbf{写表达式}：第 $k$ 个输出位 = 所有编号中该位为1的输入的OR
  \item \textbf{优先级编码器}：用 if-else 级联，优先级从高到低
  \item \textbf{常考陷阱}：$I_0$ 有效时输出000，需要GS信号区分"无输入"
\end{enumerate}
\end{examtip}

\textbf{秒杀口诀}：每个输出位 = 输入编号该位为1的所有输入的OR。

举例：
\begin{itemize}
  \item 4-2编码器：$Y_1 = I_2 + I_3$，$Y_0 = I_1 + I_3$
  \item 16-4编码器：$Y_3 = I_8 + I_9 + \cdots + I_{15}$，$Y_2 = I_4 + I_5 + I_6 + I_7 + I_{12} + I_{13} + I_{14} + I_{15}$
\end{itemize}

规律推广：对于 $2^n$-$n$ 编码器，$Y_k = \sum I_i$，其中 $i$ 的二进制表示中第 $k$ 位为1。
